My weld-seam HMI gave me a starting point in smart-machine integration: two-dimensional segmentation was connected to fixed-pose motion, while modified-DH kinematics drove an OpenGL digital twin. CR5 integration then added workspace checks, RGB-timestamp-anchored data alignment, and execution timing on real hardware. My next step is to connect such prototypes with cyber-physical production, industrial data, human interaction, and operational context. The MSc in Smart Manufacturing at The Hong Kong Polytechnic University (PolyU) is suited to that transition.

A production loop linking data to physical response would organize my study. Among subjects available under my selected route, ISE5331 Optimization and Data Analytics for Industry 4.0 turns manufacturing observations into decisions; ISE5321 Cyber-Physical Industry 4.0 Systems places those data and decisions within connected equipment; ISE5312 Industrial Human-Robot Systems and Automation extends the loop to human-aware automation; and ISE5322 Industrial Metaverse and Mixed Reality connects virtual representation with operational context. Drawing on both the HMI and CR5, I could help a team prototype a loop that makes coordinate, data, and timing assumptions explicit from digital model to physical response. That chain would advance my goal of developing intelligent-manufacturing systems whose digital and physical behavior can be evaluated together.
